% Part: incompleteness
% Chapter: arithmetization-syntax
% Section: coding-symbols

\documentclass[../../../include/open-logic-section]{subfiles}

\begin{document}

\olfileid{inc}{art}{cod}
\olsection{ترميز الرموز}

رموز اللغة الأساسية~$\Lang L$ لمنطق الرتبة الأولى هي
\[
\lfalse \quad \lnot \quad \lor \quad \land \quad \lif \quad \lforall
\quad \lexists \quad \eq \quad ( \quad ) \quad ,
\]
ومعها مجموعات !!{enumerable} من المتغيرات
وال!!{constant}s، ومجموعات !!{enumerable}
من ال!!{function}s وال!!{predicate}s على
اختلاف رتبها. ولكل رمز نعيّن \emph{رمزًا عدديًا}
واحدًا، بحيث لا يشترك رمزان مختلفان في عدد.
وإمكان ذلك معلوم من أن مجموعة الرموز
!!{enumerable}، فتوجد !!a{bijection} بينها
وبين الأعداد الطبيعية. لكننا نطلب زيادةً
على ذلك أن يُستخرج الرمز من عدده حسابيًا،
وأن تستخرج بعض خواصه، كرتبة !!a{function}.
ولهذا طرائق كثيرة؛ نختار منها طريقةً
تستعمل الدوال العودية البدائية.
وتذكر أن $\tuple{\OLId{latin-ordinary}{n}_\OLMathNumeral{0},\dots,\OLId{latin-ordinary}{n}_\OLId{latin-ordinary}{k}}$ يرمز
متتالية الأعداد $\OLId{latin-ordinary}{n}_\OLMathNumeral{0}$، و\dots، و$\OLId{latin-ordinary}{n}_\OLId{latin-ordinary}{k}$.

\begin{defn}
لكل رمز $\OLId{latin-ordinary}{s}$ من~$\Lang L$ نعرّف \emph{رمزه العددي}
$\scode \OLId{latin-ordinary}{s}$ بحسب التقسيم الآتي:
\begin{enumerate}
\item أما $\OLId{latin-ordinary}{s}$ المنطقي، فيؤخذ $\scode \OLId{latin-ordinary}{s}$ له من هذا الجدول:
\[
\begin{array}{cccccccccc}
  \lfalse & \lnot & \lor & \land & \lif & \lforall \\
\tuple{\OLMathNumeral{0}, \OLMathNumeral{0}} & \tuple{\OLMathNumeral{0}, \OLMathNumeral{1}} & \tuple{\OLMathNumeral{0}, \OLMathNumeral{2}} & \tuple{\OLMathNumeral{0}, \OLMathNumeral{3}} &
\tuple{\OLMathNumeral{0}, \OLMathNumeral{4}} & \tuple{\OLMathNumeral{0}, \OLMathNumeral{5}} \\
\lexists & \eq & ( & ) & ,\\
\tuple{\OLMathNumeral{0}, \OLMathNumeral{6}} & \tuple{\OLMathNumeral{0}, \OLMathNumeral{7}} &
\tuple{\OLMathNumeral{0}, \OLMathNumeral{8}} & \tuple{\OLMathNumeral{0}, \OLMathNumeral{9}} & \tuple{\OLMathNumeral{0}, \OLMathNumeral{10}}
\end{array}
\]
\item وأما إذا كان $\OLId{latin-ordinary}{s}$ المتغير رقم $\OLId{latin-ordinary}{i}$، أي $\Obj v_\OLId{latin-ordinary}{i}$، فنجعل
$\scode \OLId{latin-ordinary}{s}=\tuple{\OLMathNumeral{1},\OLId{latin-ordinary}{i}}$.
\item وأما إذا كان $\OLId{latin-ordinary}{s}$ ال!!{constant} رقم $\OLId{latin-ordinary}{i}$، أي~$\Obj c_\OLId{latin-ordinary}{i}$، فنجعل
  $\scode \OLId{latin-ordinary}{s}=\tuple{\OLMathNumeral{2},\OLId{latin-ordinary}{i}}$.
\item وأما إذا كان $\OLId{latin-ordinary}{s}$ ال!!{function} رقم $\OLId{latin-ordinary}{i}$ وذا الرتبة $\OLId{latin-ordinary}{n}$، أي
$\Obj f_\OLId{latin-ordinary}{i}^\OLId{latin-ordinary}{n}$، فنجعل $\scode \OLId{latin-ordinary}{s}=\tuple{\OLMathNumeral{3},\OLId{latin-ordinary}{n},\OLId{latin-ordinary}{i}}$.
\item وأما إذا كان $\OLId{latin-ordinary}{s}$ ال!!{predicate} رقم $\OLId{latin-ordinary}{i}$ وذا الرتبة $\OLId{latin-ordinary}{n}$، أي
$\Obj P_\OLId{latin-ordinary}{i}^\OLId{latin-ordinary}{n}$، فنجعل $\scode \OLId{latin-ordinary}{s}=\tuple{\OLMathNumeral{4},\OLId{latin-ordinary}{n},\OLId{latin-ordinary}{i}}$.
\end{enumerate}
\end{defn}

\begin{prop}
كل من العلاقتين الآتيتين عودية بدائية:
\begin{enumerate}
\item $\fn{Fn}(\OLId{latin-ordinary}{x},\OLId{latin-ordinary}{n})$ إذا وفقط إذا كان $\OLId{latin-ordinary}{x}$ رمز $\Obj f^\OLId{latin-ordinary}{n}_\OLId{latin-ordinary}{i}$ لبعض~$\OLId{latin-ordinary}{i}$؛
  وبعبارة أخرى، إذا كان $\OLId{latin-ordinary}{x}$ رمز !!{function} ذي رتبة $\OLId{latin-ordinary}{n}$.
\item $\fn{Pred}(\OLId{latin-ordinary}{x},\OLId{latin-ordinary}{n})$ إذا وفقط إذا كان $\OLId{latin-ordinary}{x}$ رمز $\Obj P^\OLId{latin-ordinary}{n}_\OLId{latin-ordinary}{i}$ لبعض
  $\OLId{latin-ordinary}{i}$، أو كان $\OLId{latin-ordinary}{x}$ رمز $\eq$ وكان $\OLId{latin-ordinary}{n}=\OLMathNumeral{2}$؛ وبعبارة أخرى، إذا كان $\OLId{latin-ordinary}{x}$ رمز
  !!{predicate} ذي رتبة $\OLId{latin-ordinary}{n}$.
\end{enumerate}
\end{prop}

\begin{defn}
لمتتالية الرموز $\OLId{latin-ordinary}{s}_\OLMathNumeral{0},\dots,\OLId{latin-ordinary}{s}_{\OLId{latin-ordinary}{n}-\OLMathNumeral{1}}$ نعرّف
\emph{عدد غودل} بأنه
$\tuple{\scode{\OLId{latin-ordinary}{s}_\OLMathNumeral{0}},\dots,\scode{\OLId{latin-ordinary}{s}_{\OLId{latin-ordinary}{n}-\OLMathNumeral{1}}}}$.
\end{defn}

\begin{explain}
ولا تخلط بين \emph{الرمز العددي} و\emph{عدد غودل}.
فرمز المتغير~$\Obj v_\OLMathNumeral{5}$ العددي هو
$\scode{\Obj v_\OLMathNumeral{5}}=\tuple{\OLMathNumeral{1},\OLMathNumeral{5}}=\OLMathNumeral{2}^\OLMathNumeral{2}\cdot\OLMathNumeral{3}^\OLMathNumeral{6}$.
لكن~$\Obj v_\OLMathNumeral{5}$، إذا أخذ حدًا، متتالية
رموز طولها~$\OLMathNumeral{1}$. فيكون \emph{عدد غودل}
$\Gn{\Obj v_\OLMathNumeral{5}}$ لل\emph{حد}~$\Obj v_\OLMathNumeral{5}$
هو $\tuple{\scode{\Obj v_\OLMathNumeral{5}}}=\OLMathNumeral{2}^{\scode{\Obj v_\OLMathNumeral{5}}+\OLMathNumeral{1}}
=\OLMathNumeral{2}^{\OLMathNumeral{2}^\OLMathNumeral{2}\cdot\OLMathNumeral{3}^\OLMathNumeral{6}+\OLMathNumeral{1}}$.
\end{explain}

\begin{ex}
إذا كانت $\OLId{latin-ordinary}{k}_\OLMathNumeral{0}$، و\dots، و$\OLId{latin-ordinary}{k}_{\OLId{latin-ordinary}{n}-\OLMathNumeral{1}}$ متتالية أعداد،
فرمزها $\tuple{\OLId{latin-ordinary}{k}_\OLMathNumeral{0},\dots,\OLId{latin-ordinary}{k}_{\OLId{latin-ordinary}{n}-\OLMathNumeral{1}}}$ بطريق قوى الأعداد
الأولية هو، كما تقدم،
\[
\OLMathNumeral{2}^{\OLId{latin-ordinary}{k}_\OLMathNumeral{0}+\OLMathNumeral{1}}\cdot\OLMathNumeral{3}^{\OLId{latin-ordinary}{k}_\OLMathNumeral{1}+\OLMathNumeral{1}}\cdot \dots \cdot \OLId{latin-ordinary}{p}_{\OLId{latin-ordinary}{n}-\OLMathNumeral{1}}^{\OLId{latin-ordinary}{k}_{\OLId{latin-ordinary}{n}-\OLMathNumeral{1}}+\OLMathNumeral{1}},
\]
وفيه $\OLId{latin-ordinary}{p}_\OLId{latin-ordinary}{i}$ العدد الأولي رقم $\OLId{latin-ordinary}{i}$، مع الابتداء بـ$\OLId{latin-ordinary}{p}_\OLMathNumeral{0}=\OLMathNumeral{2}$.
فالصيغة $\eq[\Obj v_\OLMathNumeral{0}][\Obj 0]$، وكتابتها الصريحة
${\eq}(\Obj v_\OLMathNumeral{0},\Obj c_\OLMathNumeral{0})$، لها عدد غودل
\[
\tuple{\scode{\eq},\scode{(},\scode{\Obj v_\OLMathNumeral{0}},\scode{,}, \scode{\Obj
    c_\OLMathNumeral{0}},\scode{)}}.
\]
فنعوّض عن $\scode{\eq}$ بـ $\tuple{\OLMathNumeral{0},\OLMathNumeral{7}}=\OLMathNumeral{2}^{\OLMathNumeral{0}+\OLMathNumeral{1}}\cdot\OLMathNumeral{3}^{\OLMathNumeral{7}+\OLMathNumeral{1}}$، وعن
$\scode{\Obj v_\OLMathNumeral{0}}$ بـ $\tuple{\OLMathNumeral{1},\OLMathNumeral{0}}=\OLMathNumeral{2}^{\OLMathNumeral{1}+\OLMathNumeral{1}}\cdot\OLMathNumeral{3}^{\OLMathNumeral{0}+\OLMathNumeral{1}}$، وهكذا.
فتكون قيمة $\Gn{=(\Obj v_\OLMathNumeral{0},\Obj c_\OLMathNumeral{0})}$
\begin{multline*}
\OLMathNumeral{2}^{\scode{=} + \OLMathNumeral{1}}\cdot \OLMathNumeral{3}^{\scode{(}+\OLMathNumeral{1}}\cdot \OLMathNumeral{5}^{\scode{\Obj v_\OLMathNumeral{0}}+\OLMathNumeral{1}}
\cdot \OLMathNumeral{7}^{\scode{,} + \OLMathNumeral{1}} \cdot \OLMathNumeral{11}^{\scode{\Obj c_\OLMathNumeral{0}}+\OLMathNumeral{1}} \cdot
\OLMathNumeral{13}^{\scode{)}+\OLMathNumeral{1}} = \\
\OLMathNumeral{2}^{\OLMathNumeral{2}^\OLMathNumeral{1}\cdot \OLMathNumeral{3}^\OLMathNumeral{8} + \OLMathNumeral{1}}\cdot \OLMathNumeral{3}^{\OLMathNumeral{2}^\OLMathNumeral{1}\cdot \OLMathNumeral{3}^\OLMathNumeral{9}+\OLMathNumeral{1}}\cdot \OLMathNumeral{5}^{\OLMathNumeral{2}^\OLMathNumeral{2}\cdot \OLMathNumeral{3}^\OLMathNumeral{1}+\OLMathNumeral{1}}
\cdot \OLMathNumeral{7}^{\OLMathNumeral{2}^\OLMathNumeral{1}\cdot \OLMathNumeral{3}^{\OLMathNumeral{11}} + \OLMathNumeral{1}} \cdot \OLMathNumeral{11}^{\OLMathNumeral{2}^\OLMathNumeral{3}\cdot\OLMathNumeral{3}^\OLMathNumeral{1}+\OLMathNumeral{1}} \cdot
\OLMathNumeral{13}^{\OLMathNumeral{2}^\OLMathNumeral{1}\cdot\OLMathNumeral{3}^{\OLMathNumeral{10}}+\OLMathNumeral{1}} = \\
\OLMathNumeral{2}^{\OLMathNumeral{13}\,\OLMathNumeral{123}}\cdot \OLMathNumeral{3}^{\OLMathNumeral{39}\,\OLMathNumeral{367}}\cdot \OLMathNumeral{5}^{\OLMathNumeral{13}}\cdot \OLMathNumeral{7}^{\OLMathNumeral{354}\,\OLMathNumeral{295}}\cdot\OLMathNumeral{11}^{\OLMathNumeral{25}}\cdot\OLMathNumeral{13}^{\OLMathNumeral{118}\,\OLMathNumeral{099}}.
\end{multline*}
\end{ex}

\end{document}
